% dtn_relation.tex
%
% Purpose:
%   Relate the Cauchy trace-subspace condition to periodic half-guide
%   Dirichlet-to-Neumann maps and limiting absorption ideas.
%
% Main algorithm:
%   No new algorithm is introduced.  The section explains when a selected
%   Cauchy subspace can be viewed as a graph N=\Lambda D.
%
% Based on:
%   notes/theory/scat_formulation.md, notes/theory/mode_selection.md, and the
%   literature on exact boundary conditions for periodic waveguides.
%
% Main changes:
%   This section separates the continuous exterior problem from the
%   finite-dimensional implementation choice.
%
% Numerical goal:
%   Make clear why the code imposes a Cauchy subspace rather than explicitly
%   forming a half-guide DtN matrix.

\section{Relation to Dirichlet-to-Neumann Conditions}

At the continuous level, a positive half-lead outgoing Dirichlet-to-Neumann map,
when it is well defined, has the form
\[
  \Lambda_{\mathrm{out}}^+:
  H^{1/2}(\Gamma_+)\to H^{-1/2}(\Gamma_+),
  \qquad
  D^+\mapsto N^+ .
\]
An analogous map may be written for the negative half-lead.  The exterior PDE
alone does not define this map.  It must be supplemented by a radiation,
admissibility, or limiting absorption condition that selects the exterior
branch.

For forced problems, uniqueness is expected only away from resonances and
eigenvalues.  For homogeneous eigenvalue problems, the nontrivial admissible
nullspace is the desired object rather than a pathology to remove.  In the exact
boundary condition framework of Joly, Li, and Fliss \cite{JolyLiFliss2006}, an
absorption parameter \(\varepsilon>0\) can be used as a regularization or
limiting absorption device to construct a unique half-guide solution and a
regularized operator \(\Lambda_\varepsilon\).  This should not be confused with
changing the final lossless eigenvalue problem; it is a way to define or
approximate the physical branch.

The implementation in this project keeps the Cauchy trace condition in subspace
form:
\[
  \begin{bmatrix}
  D^\pm\\
  N^\pm
  \end{bmatrix}
  \in
  \Eout^\pm .
\]
If the selected subspace is a graph over the Dirichlet coordinates, then one can
write
\[
  N^\pm=\Lambda_{\mathrm{out}}^\pm D^\pm .
\]
However, explicitly forming \(\Lambda_{\mathrm{out}}^\pm\) may be ill-conditioned
or impossible near points where the Dirichlet projection of the Cauchy subspace
loses rank.  The Cauchy-subspace form avoids this unnecessary graph
requirement, and it matches the data naturally produced by the Bloch
eigenproblem.
